\documentclass[11pt, a4paper]{article}
\usepackage{amsmath, amssymb, bm}
\usepackage{geometry}
\geometry{margin=2.5cm}
\setlength{\emergencystretch}{2em}

\title{Go2 Continuous Rebounce Setup and Experiment Plan}
\date{}

\begin{document}
\maketitle

\section{Current Task}

The current task is continuous rebounce control for Go2. At reset, each
environment samples a scalar target apex-height command
\begin{equation}
    h^{*} \sim \mathcal{U}(0.5, 1.2)\ \mathrm{m},
\end{equation}
and independently samples the initial drop height
\begin{equation}
    h_{\mathrm{drop}} \sim \mathcal{U}(0.5, 1.2)\ \mathrm{m}.
\end{equation}
The robot base is placed at $h_{\mathrm{drop}}$ with zero root velocity and the
default joint state. The policy must keep rebounding for the full episode. The
objective is not to jump higher forever, nor merely to recover to the initial
release height, but to make each rebound apex return close to the commanded
target height while maintaining a flat quadruped posture and staying in place.

Unlike the earlier single-rebounce setup, reaching an apex is no longer a
termination condition. The episode continues until either a failure condition is
met or the $20$ s time limit is reached. A successful rollout should therefore
show repeated cycles of
\begin{center}
    fall,\quad support interaction,\quad rebound,\quad valid apex,\quad
    re-arm after the feet return near the support.
\end{center}

\section{Environment Variants}

Two environment IDs are used:
\begin{center}
\begin{tabular}{lll}
\hline
ID & Scene & Notes \\
\hline
\texttt{Go2-Rebounce-Flat} & rigid plane & 4096 envs, env spacing $2.5$ m \\
\texttt{Go2-Rebounce-Trampoline} & deformable trampoline & 2048 envs, env spacing $4.0$ m \\
\hline
\end{tabular}
\end{center}

The trampoline variant keeps the same rebounce MDP terms but replaces the
support with a deformable trampoline and repins the trampoline rim. Task success
does not rely on rigid-deformable contact sensors. Apex and foot-clearance gates
use only root state and kinematic foot body positions.

\section{Timing}

The simulator time step and action decimation are
\begin{equation}
    \Delta t_{\mathrm{sim}} = 0.005\ \mathrm{s},
    \qquad
    d = 4,
    \qquad
    \Delta t = d\Delta t_{\mathrm{sim}} = 0.02\ \mathrm{s}.
\end{equation}
The current training episode length is
\begin{equation}
    T_{\mathrm{episode}} = 20\ \mathrm{s}.
\end{equation}
IsaacLab multiplies every weighted reward term by $\Delta t$, and logged
\texttt{Episode\_Reward/*} values are additionally divided by
$T_{\mathrm{episode}}$.

\section{Reset and Command}

The rebounce command is
\begin{equation}
    \bm{c} = [h^{*}].
\end{equation}
In the current setup, the target rebound-apex height and the initial drop height
are sampled independently from the same range:
\begin{equation}
    h^{*} \sim \mathcal{U}(0.5,1.2),
    \qquad
    h_{\mathrm{drop}} \sim \mathcal{U}(0.5,1.2).
\end{equation}
The reset event writes $h^{*}$ into the command buffer and sets the root state
to
\begin{equation}
    \bm{p}_{\mathrm{root}} =
    \bm{o}_{\mathrm{env}} +
    \left[p_x^0,\ p_y^0,\ h_{\mathrm{drop}}+h_{\mathrm{offset}}\right],
    \qquad
    \bm{v}_{\mathrm{root}}=\bm{0},
\end{equation}
where $h_{\mathrm{offset}}=0$ in the current config. Joint positions are reset to
the default Go2 pose and joint velocities are reset to zero.

During the rollout, the target command is resampled after
\begin{equation}
    t_{\mathrm{resample}} \sim \mathcal{U}(10,20)\ \mathrm{s}.
\end{equation}
This changes only $h^{*}$; it does not teleport the robot
or change the latched initial drop height $h_0$. With a $20$ s episode, most
rollouts see zero or one mid-episode target change, which is enough to test
adaptation without turning the task into rapid command following.

The command term stores the following buffers:
\begin{center}
\begin{tabular}{lll}
\hline
Symbol & Buffer & Meaning \\
\hline
$h^{*}$ & \texttt{\_target\_apex\_height} & commanded target apex height \\
$h_0$ & \texttt{drop\_height} & root height latched after reset \\
$\hat{h}_{\mathrm{apex}}$ & \texttt{last\_apex\_height} & latched valid-apex root height \\
$\hat{h}^{*}_{\mathrm{apex}}$ & \texttt{last\_apex\_target\_height} & target active at the latched valid apex \\
$b^{\mathrm{apex}}$ & \texttt{is\_apex} & one-step valid-apex flag \\
\texttt{armed} & \texttt{\_apex\_armed} & next apex can be counted \\
\hline
\end{tabular}
\end{center}

\section{Observations and Actions}

Two actor-observation variants are kept for comparison.

\paragraph{Full-state simulation observation.}
The original full-state actor observation is
\begin{equation}
    \bm{o}^{\pi,\mathrm{full}}_t =
    \left[
        h^{*},\,
        \bm{p}^{w}_t,\,
        \bm{q}^{w}_t,\,
        \bm{v}^{b}_t,\,
        \bm{\omega}^{b}_t,\,
        \tilde{\bm{q}}_t,\,
        \tilde{\dot{\bm{q}}}_t,\,
        \bm{a}_{t-1}
    \right].
\end{equation}
Here $\bm{p}^{w}_t$ is the world root position, $\bm{q}^{w}_t$ is the unique
world root quaternion, $\bm{v}^{b}_t$ is base linear velocity, and
$\bm{\omega}^{b}_t$ is base angular velocity. This variant is useful for
simulation-only baselines and for quantifying how much full-state information
helps MLP+DR.

\paragraph{Hardware-realistic deployable observation.}
The current deployable actor observation removes world root position, world
root quaternion, and base linear velocity:
\begin{equation}
    \bm{o}^{\pi,\mathrm{deploy}}_t =
    \left[
        h^{*},\,
        \bm{g}^{b}_t,\,
        \bm{\omega}^{b}_t,\,
        \tilde{\bm{q}}_t,\,
        \tilde{\dot{\bm{q}}}_t,\,
        \bm{a}_{t-1}
    \right],
\end{equation}
where $\bm{g}^{b}_t$ is the projected gravity vector in the base frame. This
observation is closer to what can be deployed without mocap: joint position,
joint velocity, IMU-like projected gravity, IMU angular velocity, previous
action, and the commanded target apex height. It intentionally does not expose
absolute base height or vertical velocity to the actor.

The deployable actor observation uses additive noise on projected gravity,
base angular velocity, joint position, and joint velocity. The critic is
asymmetric and keeps privileged full-state terms without corruption:
\[
    \bm{p}^{w}_t,\qquad
    \bm{q}^{w}_t,\qquad
    \bm{v}^{b}_t.
\]
Thus the critic can use full simulator state during training, while the actor
remains deployable. The policy action is a 12-DoF joint-position target action
with the default Go2 joint offset. In the trampoline variant, a separate action
term keeps the trampoline rim pinned.

\section{Apex Event and Metrics}

A valid apex is detected by a root vertical velocity sign change gated by the
apex re-arm state and a geometric foot-clearance condition:
\begin{equation}
    b^{\mathrm{vel}}_t =
    (\dot{z}_{t-1} > 0)
    \wedge (\dot{z}_{t} \le 0),
\end{equation}
\begin{equation}
    b^{\mathrm{feet+}}_t =
    \mathbb{I}
    \left[
        \min_{i\in\mathcal{F}} z^{\mathrm{local}}_{i,t}
        > z_{\mathrm{surface}} + 0.08
    \right],
\end{equation}
where $\mathcal{F}$ is the set of four Go2 feet and $z_{\mathrm{surface}}=0$ in
the local support frame. In the implementation this flag is named
\texttt{feet\_above\_clearance}.

To avoid counting multiple vertical velocity zero-crossings around the same
physical apex, the apex detector is armed only once per bounce. After a valid
apex is counted, it is disarmed. It is re-armed only when the feet
return below the same clearance threshold:
\begin{equation}
    b^{\mathrm{feet-}}_t =
    \mathbb{I}
    \left[
    \max_{i\in\mathcal{F}} z^{\mathrm{local}}_{i,t}
    \le z_{\mathrm{surface}} + 0.08
    \right].
\end{equation}
In the implementation this re-arm flag is named
\texttt{feet\_below\_clearance}. Thus an apex counted by the
metric/reward/timeout logic is the valid apex
\begin{equation}
    b^{\mathrm{valid}}_t =
    \texttt{armed}_t
    \wedge b^{\mathrm{vel}}_t
    \wedge b^{\mathrm{feet+}}_t.
\end{equation}

The command term is the owner of this apex state machine. When
$b^{\mathrm{valid}}_t$ fires, it latches the apex height and the target active
at that instant as \texttt{last\_apex\_height} and
\texttt{last\_apex\_target\_height}. The target is latched because command
resampling can occur after event detection and before the delayed reward reads
the event. Reward and termination terms consume this command-owned event one
manager step later; this avoids maintaining duplicate apex detectors in reward
and termination code.

The logged count metrics are:
\begin{center}
\begin{tabular}{p{0.34\linewidth}p{0.28\linewidth}p{0.26\linewidth}}
\hline
Metric & Condition & Interpretation \\
\hline
\texttt{apex\_count} & $b^{\mathrm{valid}}_t$ & valid apex count \\
\texttt{height\_matched\_apex\_count} &
$b^{\mathrm{valid}}_t \wedge |h_t-h^{*}|<0.25$ &
apex count close to target height \\
\texttt{error\_peak\_height} &
mean $|h_t-h^{*}|$ over counted apexes &
height tracking error \\
\hline
\end{tabular}
\end{center}

\section{Termination}

The current termination terms are:
\begin{center}
\begin{tabular}{p{0.35\linewidth}p{0.28\linewidth}p{0.27\linewidth}}
\hline
Term & Condition & Role \\
\hline
\texttt{base\_orientation} & base tilt angle $>0.6$ rad & failure \\
\texttt{non\_foot\_contact} & non-foot rigid contact $>1$ N & failure \\
\texttt{no\_valid\_apex\_timeout} & no valid apex for $2$ s & failure \\
\texttt{time\_out} & $20$ s reached & successful long rollout if other failures absent \\
\hline
\end{tabular}
\end{center}

\texttt{no\_valid\_apex\_timeout} reads the command-owned valid-apex pulse and
resets its timer whenever that pulse is observed. This prevents a standing
policy from simply waiting for the long time limit.

\section{Rewards}

The weighted reward is
\begin{equation}
\begin{aligned}
    R_t = \Delta t \Big(
        &-250\,r^{\mathrm{fail}}_t
        +50\,r^{\mathrm{apex}}_t \\
        &-1.5\times10^{-2}\,r^{\mathrm{energy}}_t \\
        &-2\,r^{\mathrm{orient}}_t
        -0.5\,r^{\mathrm{inplace}}_t
        -10^{-2}\,r^{\mathrm{action}}_t
        -0.05\,r^{\mathrm{joint}}_t
        -10\,r^{\mathrm{limit}}_t
    \Big).
\end{aligned}
\end{equation}

The failure pulse is
\begin{equation}
    r^{\mathrm{fail}}_t =
    \mathbb{I}[\texttt{base\_orientation}]
    + \mathbb{I}[\texttt{non\_foot\_contact}]
    + \mathbb{I}[\texttt{no\_valid\_apex\_timeout}].
\end{equation}
There is no separate normal-timeout penalty in the current setup. A rollout
that reaches \texttt{time\_out} without triggering a failure termination is
treated as a successful long rollout.

For the delayed apex pulse reward, the height error is computed from the
latched command event:
\begin{equation}
    e^{h,\mathrm{apex}}_t =
    |\texttt{last\_apex\_height}-\texttt{last\_apex\_target\_height}|.
\end{equation}
The flat-body gate is
\begin{equation}
    g^{\mathrm{flat}}_t =
    \exp\left(
        -\frac{\left\lVert \bm{g}^{b}_{t,xy}\right\rVert^2}{0.35^2}
    \right),
\end{equation}
The height reward does not include an additional soft foot-clearance gate;
foot clearance is enforced only by the valid-apex event detector.
The apex pulse reward is
\begin{equation}
    r^{\mathrm{apex}}_t =
    \mathbb{I}[b^{\mathrm{valid}}_{t-1}]
    \exp\left(-\left(\frac{e^{h,\mathrm{apex}}_t}{0.10}\right)^2\right)
    g^{\mathrm{flat}}_t.
\end{equation}
The remaining penalties are standard orientation, action-rate, joint-deviation,
and joint-position-limit regularizers. The in-place term penalizes drift from
the reset xy position and reset yaw:
\begin{equation}
    r^{\mathrm{inplace}}_t =
    \left\lVert
        \frac{\bm{p}_{xy,t}-\bm{p}_{xy,0}}{0.25}
    \right\rVert^2
    +
    \left(
        \frac{\mathrm{wrap}(\psi_t-\psi_0)}{0.5}
    \right)^2.
\end{equation}

\paragraph{Reward scale.}
In the current $20$ s setup, a single perfect apex pulse contributes
\[
    50 \cdot 0.02 / 20 = 0.05
\]
to the logged \texttt{Episode\_Reward/rebounce\_height}. A single failure
termination contributes approximately
\[
    -250 \cdot 0.02 / 20 = -0.25
\]
for the environments where that failure happened, before averaging over all
reset environments. With the current absolute-work weight, a typical sparse
energy pulse of $W^{\mathrm{abs}}/h^{*}\approx 250$ contributes
\[
    -1.5\times10^{-2}\cdot250\cdot0.02/20 \approx -0.00375
\]
to the logged \texttt{Episode\_Reward/energy\_penalty} for one valid apex.

\section{Implementation Files}

\begin{center}
\begin{tabular}{p{0.42\linewidth}p{0.46\linewidth}}
\hline
File & Purpose \\
\hline
\texttt{go2\_rebounce\_env\_cfg.py} & scene, reset, reward, termination config \\
\texttt{mdp/commands.py} & rebounce command, apex metrics, re-arm state, energy metrics \\
\texttt{mdp/events.py} & reset from sampled drop height \\
\texttt{mdp/rewards.py} & height reward, energy penalty, and failure pulses \\
\texttt{mdp/terminations.py} & timeout-without-apex termination \\
\texttt{scripts/rsl\_rl/play-rebounce.py} & play script with apex height and energy printout \\
\hline
\end{tabular}
\end{center}

\section{Experiment Roadmap}

The experiment order should separate task-definition changes from style,
efficiency, and final ablations.

\subsection{Step 0: Freeze the Current Continuous-Rebounce Baseline}

First train and save a clean baseline with the current config. Record the git
state, run ID, checkpoint, and videos. The expected healthy indicators are:
\begin{itemize}
    \item \texttt{Episode\_Termination/time\_out} close to $1$,
    \item \texttt{no\_valid\_apex\_timeout} close to $0$,
    \item \texttt{apex\_count} close to the visually observed bounce count,
    \item \texttt{height\_matched\_apex\_count}/\texttt{apex\_count} high,
    \item \texttt{error\_peak\_height} low.
\end{itemize}

\subsection{Step 1: Decouple Initial Drop Height from Target Apex Height}

This has been implemented in the current setup. The target height and the
initial release height are no longer the same sampled scalar:
\[
    h^{*} \sim \mathcal{U}(0.5,0.8), \qquad
    h_{\mathrm{drop}} \sim \mathcal{U}(0.5,0.8).
\]
The height rewards and metrics now compare apex height against $h^{*}$ rather
than against $h_{\mathrm{drop}}$. This tests whether the policy controls
rebound height, rather than only recovering to the height from which it was
dropped.

\subsection{Step 2: Add In-Place Hopping}

This has also been implemented in the current setup with a small dense penalty
on displacement from the reset xy position and reset yaw. Suggested diagnostics
for the run are
\[
    \lVert \bm{p}_{xy,t}-\bm{p}_{xy,0}\rVert,\qquad
    \lVert \bm{v}_{xy,t}\rVert,\qquad
    |\psi_t-\psi_0|,\qquad
    |\dot{\psi}_t|.
\]
The current weight is intentionally modest so vertical rebounce quality is not
suppressed. The goal is to keep the robot rebounding in place without
sacrificing apex-height tracking.

\subsection{Step 3: Minimal Leg-Motion Regulation}

Before optimizing energy, first remove obvious high-frequency or visually
unacceptable leg motion. This should be a light regulation pass, not a strong
posture-imitation objective. The goal is to prevent wasted actuator motion and
unsafe leg swings while still allowing the robot to shape its body and legs to
use the trampoline.

Test the following regularization variants one at a time:
\begin{enumerate}
    \item increase \texttt{joint\_deviation\_l1} weight, e.g.
    $-0.05 \rightarrow -0.10$ first; avoid jumping immediately to a strong
    posture penalty because it can suppress useful trampoline-exploitation
    strategies;
    \item add a small joint-velocity penalty to reduce fast flight leg motion;
    \item test whether the command-side hard foot-clearance apex gate is still
    needed after removing the reward-side soft foot-clearance gate.
\end{enumerate}
This step should be done before energy optimization so that the energy term is
not merely compensating for a leg-motion artifact. After each variant, check
\texttt{apex\_count}, \texttt{height\_matched\_apex\_count},
\texttt{error\_peak\_height}, in-place drift, and videos. If height tracking or
survival degrades, the regularization is too strong.

\subsection{Step 4: Add Energy Optimization on a Fixed Trampoline}

Only after the behavior is stable and visually acceptable, add energy terms on
the fixed trampoline. Keeping trampoline parameters fixed during this step
isolates the question: can the policy maintain commanded apex height with less
active work by using the trampoline's passive energy storage?

Two mechanical-work definitions should be tested:
\[
    P^{+}_t =
    \sum_j \max(\tau_{j,t}\dot{q}_{j,t}, 0),
    \qquad
    P^{\mathrm{abs}}_t =
    \sum_j |\tau_{j,t}\dot{q}_{j,t}|.
\]
The reward is sparse. Between two valid apex events, the energy command
accumulates
\[
    W^{+}_k=\sum_{t\in k}P^{+}_t\Delta t,
    \qquad
    W^{\mathrm{abs}}_k=\sum_{t\in k}P^{\mathrm{abs}}_t\Delta t.
\]
At the next valid apex, the reward penalty is normalized by the commanded target
height $h^{*}_k$:
\[
    c^{\mathrm{energy}}_k =
    \frac{W_k}{\max(h^{*}_k,\epsilon)}.
\]
Using the target height for the reward prevents the policy from jumping higher
than requested only to dilute the work-per-height penalty. The reported metrics
below still use the achieved apex height, which keeps them interpretable as the
actual work needed per meter of realized rebound height.

The implementation exposes this through the \texttt{mode} parameter of
\texttt{energy\_penalty} in \texttt{go2\_rebounce\_env\_cfg.py}. Set
\texttt{mode=positive} for $W^{+}_k/h^{*}_k$ or
\texttt{mode=absolute} for $W^{\mathrm{abs}}_k/h^{*}_k$. The raw reward follows
the same convention as the apex-height pulse: Isaac Lab applies the global
$\Delta t$ scaling, and the event scale is handled in the configured weight.
With the current $\Delta t=0.02$ s, the current absolute-work trial weight is
\[
    w_{\mathrm{energy}} = -1.5\times10^{-2}.
\]
The positive-work term penalizes active mechanical energy injected by the
motors:
\[
    \tau_{j,t}\dot{q}_{j,t} > 0.
\]
It is the preferred first test because it discourages the robot from actively
``kicking'' to create apex height while still allowing negative work for
landing stabilization. The absolute-work term penalizes both positive and
negative mechanical work, so it also discourages active braking:
\[
    r^{\mathrm{abs\_work}}_t
    =
    r^{\mathrm{pos\_work}}_t
    +
    \sum_j \max(-\tau_{j,t}\dot{q}_{j,t},0).
\]
It should be tested after the positive-work variant, or with a much smaller
weight, because too much absolute-work penalty can reduce useful damping and
postural stabilization during trampoline compression.

A torque-squared term,
\[
    \sum_j \tau_{j,t}^{2},
\]
is also useful as a regularizer, but it mainly discourages large torques rather
than directly measuring active mechanical work.

The previous positive-work run shows the intended first-order behavior: the
logged \texttt{Episode\_Reward/energy\_penalty} is roughly $10\%$ of
\texttt{Episode\_Reward/rebounce\_height}, the valid apex count increases, and
\texttt{positive\_work\_per\_height} decreases. This indicates that the energy
term is influencing the policy without dominating the height-tracking task.

Use a small weight or a curriculum that turns on the energy penalty only after
the policy already achieves stable rebounce. Keep the height and survival terms
dominant. If energy decreases while \texttt{apex\_count},
\texttt{height\_matched\_apex\_count}, \texttt{error\_peak\_height}, or
\texttt{time\_out} degrades, the energy weight is too large or was introduced
too early. Compare policies with
positive-work-per-height and absolute-work-per-height; use braking ratio as a
diagnostic. A low positive-work policy with high braking ratio may simply be
absorbing trampoline energy through active braking rather than using the
trampoline efficiently.

The energy command term logs:
\begin{center}
\scriptsize
\begin{tabular}{p{0.42\linewidth}p{0.46\linewidth}}
\hline
Metric & Definition \\
\hline
\texttt{Metrics/energy/positive\_work\_per\_height} &
mean over valid apexes of $W_k^{+}/h_k^{\mathrm{apex}}$ \\
\texttt{Metrics/energy/absolute\_work\_per\_height} &
mean over valid apexes of $W_k^{\mathrm{abs}}/h_k^{\mathrm{apex}}$ \\
\texttt{Metrics/energy/braking\_ratio} &
$W^{-}_{\mathrm{episode}}/\max(W^{\mathrm{abs}}_{\mathrm{episode}},\epsilon)$ \\
\hline
\end{tabular}
\end{center}

\subsection{Step 5: Add Trampoline Parameter Randomization}

After the energy-optimized policy works on the fixed trampoline, introduce
trampoline randomization. The current Go2 rebounce trampoline setup uses
multi-parameter randomization over stiffness, mass, friction, material damping,
and Poisson ratio:
\[
    \log E \sim
    \mathcal{U}\left(\log(2.0\times 10^4),\log(8.0\times 10^4)\right),
    \qquad
    m \sim \mathcal{U}(5,15),
\]
\[
    \mu_d \sim \mathcal{U}(0.4,1.2),
    \qquad
    d_e \sim \mathcal{U}(0.01,0.1),
    \qquad
    \nu \sim \mathcal{U}(0.25,0.45),
\]
with \texttt{damping\_scale}=1.0 fixed. Young's modulus is sampled
log-uniformly because stiffness changes multiplicatively; the other parameters
are sampled uniformly in their current engineering ranges. The earlier
first-stage setup used
\[
    E\sim \mathrm{LogUniform}(4.0\times10^4,1.6\times10^5),
    \qquad m=10,
\]
and remains a useful historical baseline. The current range intentionally
shifts toward softer trampolines because the old $8.0\times10^4$ and
$1.6\times10^5$ settings were too similar in height tracking.

The randomizable trampoline parameters should be introduced in stages:
\begin{center}
\small
\begin{tabular}{p{0.24\linewidth}p{0.38\linewidth}p{0.24\linewidth}}
\hline
Parameter & Meaning & Suggested use \\
\hline
$E$ / \texttt{youngs\_modulus} & membrane stiffness; controls soft versus
stiff rebound timing & first-stage randomization \\
$m$ / \texttt{mass} & trampoline inertia; changes response speed and energy
exchange timing & first-stage randomization \\
\texttt{elasticity\_damping} or \texttt{damping\_scale} & material energy
dissipation & second-stage, one damping knob at a time \\
\texttt{dynamic\_friction} & tangential foot--trampoline friction, affecting
slip and in-place stability & robustness sweep \\
$\nu$ / \texttt{poissons\_ratio} & deformation mode and volume preservation &
late-stage small-range sweep \\
\hline
\end{tabular}
\end{center}

Keep numerical and sleep parameters fixed during the main physics
randomization study: \texttt{vertex\_velocity\_damping},
\texttt{sleep\_damping}, \texttt{sleep\_threshold},
\texttt{settling\_threshold}, solver iteration counts, mesh resolution,
\texttt{contact\_offset}, and \texttt{rest\_offset}. These parameters can affect
stability and numerical dissipation, but they are less direct descriptions of a
real trampoline than stiffness, mass, damping, friction, and Poisson ratio.
For visual parameter sweeps, the rebounce play script exposes fixed overrides
for $E$, $m$, \texttt{dynamic\_friction},
\texttt{elasticity\_damping}, \texttt{damping\_scale}, and
\texttt{poissons\_ratio}.

Use the stiffness range as the first-stage domain-randomization test:
\[
    \log E \sim \mathcal{U}(\log(2.0\times10^4),\log(8.0\times10^4)).
\]
If this range is too difficult, narrow it temporarily, for example to
$E/E_0\in[0.8,1.25]$, to separate learning instability from adaptation
failure. If stiffness randomization is stable, add mass and material damping
next, then friction and Poisson ratio. Randomizing too many parameters too
early can make it difficult to tell whether a behavior change came from energy
optimization, leg regulation, or trampoline dynamics. It can also encourage a
more conservative active-jumping strategy that is robust but less clearly
trampoline-efficient.

The randomized parameters should not be provided directly to the deployable
actor policy, since the real robot does not know the true trampoline stiffness,
mass, or damping. For the paper narrative, the adaptation methods should be
tested in increasing order of complexity:
\begin{center}
\small
\begin{tabular}{p{0.20\linewidth}p{0.42\linewidth}p{0.26\linewidth}}
\hline
Method & Description & Purpose \\
\hline
MLP + DR & train one memoryless policy across randomized trampoline stiffness;
no trampoline parameters and no history in the actor observation & baseline
robust policy \\
MLP + history & stack recent observations/actions and feed the flattened history
to an MLP & test whether short-term interaction history is enough \\
RNN policy & use an LSTM/GRU policy to infer hidden trampoline dynamics from
longer temporal context & test recurrent online adaptation \\
RMA / latent estimator & train with privileged dynamics information, then learn
a deployable latent estimator from history & final sim-to-real adaptation story
\\
\hline
\end{tabular}
\end{center}
Privileged trampoline parameters may be used by the critic, teacher, or
adaptation module during training, but not as direct actor inputs for a policy
intended to transfer to hardware.

The first method is intentionally a partially observable baseline. Since the
actor sees neither the trampoline parameter nor interaction history, it cannot
condition its timing or leg impedance on the current stiffness. The expected
failure mode is therefore an averaged strategy: performance on the soft
trampoline, e.g. $E=4.0\times10^4$, may improve compared with a policy trained
only at $E=8.0\times10^4$, but nominal-stiffness performance may degrade and
the hard trampoline, e.g. $E=1.6\times10^5$, may require a different timing
compromise. This creates the justification for the later methods: history,
recurrence, or an RMA-style latent estimator should recover stiffness-dependent
behavior while keeping the deployable actor free of privileged trampoline
parameters.

\paragraph{Preliminary fixed-condition result.}
A first evaluation sweep compared two policies: a stiffness-randomized policy
trained with
\[
    E\sim \mathrm{LogUniform}(4.0\times10^4,1.6\times10^5)
\]
and a fixed-trampoline policy trained only at
$E=8.0\times10^4$, $m=10$ kg. Both policies were evaluated on fixed
conditions
\[
    E\in\{4.0\times10^4,8.0\times10^4,1.6\times10^5\},
    \qquad
    h^*\in\{0.5,0.65,0.8\}.
\]
The stiffness-randomized MLP achieved $100\%$ success on all nine conditions
and maintained low height error, with overall
\[
    \mathrm{MAE}=0.014\ \mathrm{m},
    \qquad
    |\mathrm{bias}|=0.011\ \mathrm{m}.
\]
The fixed-$8.0\times10^4$ policy performed well at and above its training
stiffness, but failed to generalize to the soft trampoline. In particular, at
$(E,h^*)=(4.0\times10^4,0.5)$ it had $0\%$ success due to
\texttt{no\_valid\_apex\_timeout}, and at
$(E,h^*)=(4.0\times10^4,0.65)$ it systematically under-jumped with
\[
    \mathrm{MAE}=0.103\ \mathrm{m},
    \qquad
    h/h^*=0.842.
\]
This supports the first-stage story that domain randomization improves
robustness to trampoline stiffness. However, the randomized policy is not
strictly better in all metrics: compared with the fixed policy, it uses more
motor work per target height and has higher orientation RMS on the nominal and
stiff trampolines. Averaged over all fixed evaluation settings, the randomized
policy had
\[
    W^{+}/h^* \approx 159,
    \quad
    W^{\mathrm{abs}}/h^* \approx 247,
    \quad
    \mathrm{orientation\_rms}\approx0.092,
\]
while the fixed policy had
\[
    W^{+}/h^* \approx 138,
    \quad
    W^{\mathrm{abs}}/h^* \approx 207,
    \quad
    \mathrm{orientation\_rms}\approx0.063.
\]
Thus the current MLP+DR policy is a robust average strategy: it solves height
tracking across stiffness, but it appears less energy-efficient and slightly
less posture-stable than a policy specialized to the nominal trampoline.

\paragraph{Latest trampoline-parameter sweep.}
A later fixed-condition sweep evaluated the current DR policy on $27$
conditions:
\[
    h^*\in\{0.5,0.85,1.2\},
    \qquad
    E\in\{2.0\times10^4,4.0\times10^4,8.0\times10^4\},
    \qquad
    d_e\in\{0.01,0.05,0.1\},
\]
with $m=10$, $\mu_d=0.75$, \texttt{damping\_scale}=1.0, and
$\nu=0.35$. Each condition used $64$ rollouts. The policy was stable:
$26/27$ conditions had $100\%$ success, and the only non-perfect condition was
$(h^*,E,d_e)=(1.2,8.0\times10^4,0.01)$ with one
\texttt{no\_valid\_apex\_timeout} failure.

Height tracking remains strong for moderate targets but the high target becomes
a useful stress test:
\[
\begin{array}{c|ccc}
h^* & 0.5 & 0.85 & 1.2 \\
\hline
\mathrm{MAE}\ [\mathrm{m}] & 0.016 & 0.024 & 0.140 \\
\mathrm{bias}\ [\mathrm{m}] & +0.004 & -0.009 & -0.115 \\
h/h^* & 1.009 & 0.989 & 0.904
\end{array}
\]
Thus the main failure mode at $h^*=1.2$ is not random falling but systematic
under-jumping. Damping has the clearest effect:
\[
\begin{array}{c|ccc}
d_e & 0.01 & 0.05 & 0.10 \\
\hline
\mathrm{MAE}\ [\mathrm{m}] & 0.035 & 0.045 & 0.101 \\
\mathrm{bias}\ [\mathrm{m}] & +0.011 & -0.035 & -0.096 \\
W^+/h^* & 211 & 281 & 309
\end{array}
\]
Low damping is easier and more energy-efficient but produces more xy drift;
high damping dissipates trampoline energy, increases active motor work, and
causes under-jumping, especially for $h^*=1.2$. This suggests that future
adaptation methods should be evaluated on high-target/high-damping stress
settings in addition to ordinary moderate-height tracking.

\paragraph{Observation-ablation hypothesis for RNN/RMA.}
The earlier actor observed root position and base linear velocity. Those terms
make the instantaneous observation highly informative: root height and vertical
velocity expose bounce phase, impact timing, and part of the trampoline
response. The current actor removes those terms and keeps only the deployable
observation set, while the critic remains privileged. This makes the task more
realistic for hardware without mocap and creates a better test bed for online
dynamics inference. The next question is whether a memoryless MLP can still
infer enough from the instantaneous deployable observation, or whether history,
recurrence, or an RMA latent are needed.

The proposed controlled variants are:
\begin{center}
\small
\begin{tabular}{p{0.24\linewidth}p{0.42\linewidth}p{0.24\linewidth}}
\hline
Variant & Actor observation & Purpose \\
\hline
Historical full-observation MLP+DR & includes root position and base linear
velocity & upper-bound memoryless baseline \\
Deployable MLP+DR & no root position, no root quaternion, no base linear
velocity & current memoryless baseline \\
Deployable history MLP & same deployable observation, with stacked
history/actions & test short-horizon implicit velocity/dynamics inference \\
Deployable RNN/RMA & same deployable observation, with recurrent state or
learned latent dynamics & test online trampoline
identification \\
\hline
\end{tabular}
\end{center}

This ablation should be framed carefully. It should not artificially remove
sensors only to make the baseline fail. It is justified if real deployment does
not provide reliable base position or base linear velocity, or if mocap/state
estimator signals are noisy and delayed. In that case, removing or heavily
corrupting these observations is a realistic sim-to-real constraint. The fair
comparison is then not ``full-observation MLP versus reduced-observation
RNN'', but rather MLP, history-MLP, RNN, and RMA under the same deployable
observation set. If reduced-observation MLP+DR loses height tracking while
RNN/RMA recovers it, the paper can claim that temporal inference or latent
dynamics estimation is needed for robust trampoline adaptation. If height
tracking remains good even after the reduction, then RNN/RMA should be judged
mainly on energy efficiency, posture stability, and adaptation speed.

\paragraph{Next adaptation sequence.}
The next method-development sequence is:
\begin{enumerate}
    \item \textbf{MLP + observation history}: keep the deployable instantaneous
    observation, stack a short window of recent observations/actions, and feed
    the flattened vector to the same MLP policy. This is the simplest test of
    whether recent interaction history reveals trampoline damping and timing.
    \item \textbf{RNN policy}: replace the flattened history with a recurrent
    policy state, e.g. GRU or LSTM. This tests whether learned memory improves
    online adaptation over a fixed-length history window.
    \item \textbf{RMA}: train a teacher or privileged module with trampoline
    dynamics information, then train a deployable adaptation module that
    infers a latent dynamics code from history. This is the final paper-facing
    adaptation story if MLP+history or RNN expose a clear benefit.
\end{enumerate}
All three methods should use the same deployable observation set for a fair
comparison. The privileged trampoline parameters may be used for critic,
teacher, or latent-supervision losses, but not as direct actor inputs at test
time.

To make different trampoline dynamics visible in evaluation, run fixed
evaluation sweeps over several discrete trampoline settings instead of relying
only on random training averages. For each setting, compare height error,
apex count, time-out rate, motor positive work, torque norm, contact/stance
duration, and videos. A useful result should show that the policy adapts its
timing, leg motion, and motor work across soft and stiff trampolines while
maintaining the same commanded apex-height objective.

For paper-level comparison, evaluate each method on fixed conditions rather
than only on the training distribution. A condition is a pair such as
$(E,h^*)=(4.0\times10^4,0.5)$, where the trampoline stiffness and target apex
height are held fixed for each $20$ s rollout. Run $N$ independent rollouts per
condition. Compute metrics within each rollout first, then report the mean and
standard deviation across rollouts, so that an episode with many apexes does
not dominate the statistics.

\begin{center}
\scriptsize
\begin{tabular}{p{0.32\linewidth}p{0.43\linewidth}p{0.15\linewidth}}
\hline
Metric & Definition & Purpose \\
\hline
\texttt{success\_rate} &
fraction of rollouts that reach the $20$ s timeout without
\texttt{base\_orientation}, \texttt{non\_foot\_contact}, or
\texttt{no\_valid\_apex\_timeout} & stability \\
\texttt{failure\_distribution} &
counts or fractions of each failed-termination reason & failure diagnosis \\
\texttt{apex\_count\_per\_20s} &
number of valid apexes in one rollout & sustained hopping frequency \\
\texttt{height\_matched\_apex\_count} &
number of valid apexes satisfying the height tolerance & useful successful hops \\
\texttt{height\_mae} &
$K^{-1}\sum_{k=1}^{K}|h_k-h^*_k|$ over valid apexes in the rollout &
absolute tracking accuracy \\
\texttt{height\_rmse} &
$\sqrt{K^{-1}\sum_{k=1}^{K}(h_k-h^*_k)^2}$ & sensitivity to large errors \\
\texttt{height\_bias} &
$K^{-1}\sum_{k=1}^{K}(h_k-h^*_k)$ & systematic under/over-shoot \\
\texttt{mean\_h\_over\_target} &
$K^{-1}\sum_{k=1}^{K}h_k/h^*_k$ & normalized achieved height \\
\texttt{height\_p90\_error} &
$90$th percentile of $|h_k-h^*_k|$ & tail robustness \\
\texttt{pos\_work\_per\_target\_height} &
mean over valid apex intervals of $W_k^+/h^*_k$ & active work cost \\
\texttt{abs\_work\_per\_target\_height} &
mean over valid apex intervals of $W_k^{\mathrm{abs}}/h^*_k$ &
total motor work cost \\
\texttt{braking\_ratio} &
$W^-_{\mathrm{episode}}/\max(W^{\mathrm{abs}}_{\mathrm{episode}},\epsilon)$ &
braking versus propulsion balance \\
\texttt{xy\_drift}, \texttt{yaw\_drift} &
final or RMS base displacement/yaw error over the rollout & in-place behavior \\
\texttt{orientation\_rms} &
RMS roll/pitch or projected-gravity error & posture quality \\
\hline
\end{tabular}
\end{center}

For height metrics, use the valid apexes detected by the command term. If
$K=0$, the rollout should be counted as a failure and the height metrics should
be excluded from the successful-rollout mean or reported as undefined. For
energy metrics, the target-height denominator is preferred for method
comparison because it does not reward jumping higher than commanded; achieved
height normalization can still be logged as an efficiency diagnostic.

\subsection{TODO: Real-Robot Observation Noise Measurement}

Before finalizing sim-to-real observation corruption, measure the observation
noise on the physical robot and use the measured ranges to set the policy
observation noise. The required measurements are:
\begin{itemize}
    \item joint position noise from the robot encoder readings;
    \item joint velocity noise from the robot state estimator or finite
    differences of encoder readings;
    \item IMU angular velocity noise for the base angular-velocity observation;
    \item mocap or state-estimator base position noise;
    \item mocap or state-estimator base linear-velocity noise.
\end{itemize}

For each signal, record at least a static standing segment and a representative
rebounce or hopping segment. Estimate bias, standard deviation, outliers,
latency, and any filtering delay. The resulting noise magnitudes should be used
to update the observation corruption in \texttt{go2\_rebounce\_env\_cfg.py},
especially for joint position, joint velocity, base angular velocity, base
position, and base linear velocity.

\subsection{Step 6: Final Ablation Experiments}

Run formal ablations after the final task definition is stable. Suggested
variants are:
\begin{center}
\begin{tabular}{p{0.22\linewidth}p{0.34\linewidth}p{0.34\linewidth}}
\hline
Variant & Change & Main question \\
\hline
Full & final setup & reference \\
No apex timeout & remove \texttt{no\_valid\_apex\_timeout} & is the anti-wait gate needed? \\
No foot gate & remove the command-side valid-apex foot-clearance gate & does it prevent standing or foot exploits? \\
No in-place reward & remove $xy/yaw$ terms & does the policy drift? \\
No regulation & remove flight-leg regulation & does leg motion return? \\
No energy reward & remove energy term & efficiency versus performance \\
No trampoline DR & fix trampoline parameters & robustness versus specialization \\
\hline
\end{tabular}
\end{center}

For every variant, compare success rate, failure distribution, height error,
apex count, height-matched apex count, drift metrics, energy, trampoline
setting, and videos.

\section{Recommended Order}

The recommended order is:
\begin{enumerate}
    \item keep the current deployable-observation MLP+DR baseline as the
    reference;
    \item add observation history to the MLP and evaluate the same fixed
    trampoline conditions;
    \item test RNN policies under the same observation and evaluation protocol;
    \item move to RMA only after the history/RNN baselines reveal which
    trampoline conditions need online adaptation;
    \item measure real-robot observation noise and update observation
    corruption;
    \item run final ablations over reward, termination, energy, and trampoline
    randomization terms.
\end{enumerate}

This order treats MLP+DR as the current baseline, then adds temporal inference
in increasing complexity: explicit history, recurrent memory, and finally an
RMA-style latent dynamics estimator.

\end{document}
